\documentclass[12pt,letterpaper]{article}

\usepackage[T1]{fontenc}
\usepackage[spanish,es-tabla]{babel}
\usepackage{helvet}
\renewcommand{\familydefault}{\sfdefault}
\usepackage{geometry}
\geometry{left=2.5cm,right=2.5cm,top=2.5cm,bottom=2.5cm}
\usepackage{amsmath,amssymb}
\usepackage{graphicx}
\usepackage{float}
\usepackage{booktabs}
\usepackage{array}
\usepackage{xcolor}
\usepackage{listings}
\usepackage{setspace}
\usepackage{ragged2e}
\usepackage{hyperref}
\usepackage{caption}
\usepackage{longtable}

\hypersetup{
    colorlinks=true,
    linkcolor=black,
    urlcolor=blue,
    citecolor=black
}

\setstretch{1.15}
\justifying

\definecolor{codegray}{RGB}{245,245,245}
\definecolor{codeborder}{RGB}{180,180,180}
\definecolor{codeblue}{RGB}{35,65,130}
\definecolor{codegreen}{RGB}{35,120,60}
\definecolor{codered}{RGB}{150,45,45}

\lstdefinestyle{pythonstyle}{
    language=Python,
    basicstyle=\ttfamily\small,
    backgroundcolor=\color{codegray},
    frame=single,
    rulecolor=\color{codeborder},
    breaklines=true,
    showstringspaces=false,
    tabsize=4,
    numbers=left,
    numberstyle=\tiny\color{black!60},
    keywordstyle=\color{codeblue}\bfseries,
    commentstyle=\color{codegreen},
    stringstyle=\color{codered},
    columns=fullflexible,
    keepspaces=true
}

\begin{document}

\begin{titlepage}
\centering

\IfFileExists{UQ.jpg}{
    \includegraphics[width=3.2cm]{UQ.jpg}\par\vspace{0.5cm}
}{
    \vspace*{1cm}
}

{\Large \textbf{Universidad del Quind\'{i}o}}\par
\vspace{0.2cm}
{\large Facultad de Ciencias B\'{a}sicas y Tecnolog\'{i}as}\par
\vspace{0.2cm}
{\large Programa de F\'{i}sica}\par

\vspace{2.2cm}

{\Large \textbf{Proyecto final de Programaci\'{o}n}}\par
\vspace{0.6cm}
{\LARGE \textbf{Simulaci\'{o}n y optimizaci\'{o}n de un lanzamiento parab\'{o}lico hacia un objetivo}}\par

\vspace{2.2cm}

\begin{flushleft}
\textbf{Estudiante:} Juan David C\'{a}rdenas V\'{e}lez\\
\textbf{Asignatura:} Programaci\'{o}n\\
\textbf{Profesor:} Santiago Echeverry\\
\textbf{Fecha:} \today
\end{flushleft}

\vfill
{\large Armenia, Quind\'{i}o}\par
\end{titlepage}

\tableofcontents
\newpage

\section{Introducci\'{o}n}

En este proyecto se desarroll\'{o} una simulaci\'{o}n computacional de un lanzamiento parab\'{o}lico. La idea principal fue construir un programa que, a partir de una velocidad inicial, una gravedad y la posici\'{o}n de un objetivo, encuentre el \'{a}ngulo de lanzamiento que hace que el proyectil pase lo m\'{a}s cerca posible de ese punto.

El trabajo mezcla una parte de f\'{i}sica bastante conocida, que es el movimiento parab\'{o}lico, con una parte de programaci\'{o}n orientada a objetos. La intenci\'{o}n no fue hacer un c\'{o}digo demasiado complicado, sino un programa que se pueda entender y defender bien. Por eso el modelo f\'{i}sico se mantuvo ideal: no se considera rozamiento del aire, la gravedad es constante y el proyectil parte desde el origen.

El programa utiliza \texttt{NumPy} para los c\'{a}lculos num\'{e}ricos, \texttt{SciPy} para buscar el \'{a}ngulo que mejor se ajusta al objetivo y \texttt{Matplotlib} para mostrar la trayectoria y animar el movimiento. Adem\'{a}s, el c\'{o}digo se organiz\'{o} con clases, herencia y una clase abstracta, de forma que tambi\'{e}n cumpliera con la parte de dise\~{n}o computacional pedida en el proyecto.

\section{Planteamiento del problema}

El problema que se quiere resolver es el siguiente: se tiene un proyectil que se lanza desde el origen con una velocidad inicial conocida. Tambi\'{e}n se tiene un objetivo ubicado en una posici\'{o}n determinada del plano, dada por las coordenadas \((x_{\text{objetivo}}, y_{\text{objetivo}})\). La pregunta es qu\'{e} \'{a}ngulo de lanzamiento permite que la trayectoria del proyectil pase lo m\'{a}s cerca posible de ese objetivo.

En mec\'{a}nica ideal, si se conocen las condiciones iniciales, la trayectoria queda definida. Sin embargo, desde el punto de vista computacional, el programa necesita buscar entre varios \'{a}ngulos posibles para decidir cu\'{a}l se ajusta mejor al objetivo. Por eso se usa una funci\'{o}n de optimizaci\'{o}n: no porque haya error experimental, sino porque el computador necesita una medida matem\'{a}tica para comparar \'{a}ngulos.

La medida usada en el programa es la diferencia vertical entre la altura del proyectil y la altura del objetivo cuando el proyectil llega a la misma posici\'{o}n horizontal del blanco. En el c\'{o}digo aparece como un ``error'', pero se debe entender m\'{a}s bien como una desviaci\'{o}n vertical respecto al objetivo.

\section{Objetivos}

\subsection{Objetivo general}

Desarrollar una simulaci\'{o}n computacional del movimiento parab\'{o}lico de un proyectil, usando programaci\'{o}n orientada a objetos, para encontrar el \'{a}ngulo de lanzamiento que permite acercarse a un objetivo definido en el plano.

\subsection{Objetivos espec\'{i}ficos}

\begin{itemize}
    \item Implementar las ecuaciones del movimiento parab\'{o}lico en Python.
    \item Usar \texttt{NumPy} para calcular las posiciones del proyectil en diferentes instantes de tiempo.
    \item Usar \texttt{SciPy} para optimizar el \'{a}ngulo de lanzamiento.
    \item Visualizar la trayectoria mediante \texttt{Matplotlib}.
    \item Animar el movimiento del proyectil usando \texttt{FuncAnimation}.
    \item Organizar el c\'{o}digo mediante clases, herencia y una clase abstracta.
\end{itemize}

\section{Marco te\'{o}rico}

El movimiento parab\'{o}lico se puede ver como la combinaci\'{o}n de dos movimientos: uno horizontal y uno vertical. En este proyecto se toma el caso ideal donde no hay rozamiento del aire. Por eso, en el eje horizontal no hay aceleraci\'{o}n y la velocidad horizontal permanece constante.

La posici\'{o}n horizontal se calcula con

\begin{equation}
x(t)=v_0\cos(\theta)t,
\end{equation}

donde \(v_0\) es la velocidad inicial, \(\theta\) es el \'{a}ngulo de lanzamiento y \(t\) es el tiempo.

En el eje vertical s\'{i} act\'{u}a la gravedad. Como se toma hacia arriba como direcci\'{o}n positiva, la aceleraci\'{o}n vertical es \(-g\). Por eso la posici\'{o}n vertical est\'{a} dada por

\begin{equation}
y(t)=v_0\sin(\theta)t-\frac{1}{2}gt^2.
\end{equation}

El primer t\'{e}rmino, \(v_0\sin(\theta)t\), representa la subida inicial del proyectil debido a la componente vertical de la velocidad. El segundo t\'{e}rmino, \(-\frac{1}{2}gt^2\), representa el efecto de la gravedad, que va curvando la trayectoria hacia abajo.

Para simular la trayectoria completa se calcula el tiempo de vuelo. Como el proyectil parte desde \(y=0\) y se simula hasta que vuelve al suelo, se usa

\begin{equation}
t_f=\frac{2v_0\sin(\theta)}{g}.
\end{equation}

Esta expresi\'{o}n sale de imponer \(y(t)=0\) en la ecuaci\'{o}n vertical y tomar la soluci\'{o}n distinta de \(t=0\).

La optimizaci\'{o}n entra cuando se quiere encontrar el \'{a}ngulo m\'{a}s conveniente. Para un \'{a}ngulo dado, primero se calcula el tiempo en que el proyectil llega horizontalmente al objetivo. A partir de la ecuaci\'{o}n horizontal,

\begin{equation}
x=v_0\cos(\theta)t,
\end{equation}

se despeja

\begin{equation}
t=\frac{x_{\text{objetivo}}}{v_0\cos(\theta)}.
\end{equation}

Luego se eval\'{u}a la altura del proyectil en ese instante y se compara con la altura del objetivo. La cantidad que el programa minimiza es

\begin{equation}
\left|y_{\text{proyectil}}-y_{\text{objetivo}}\right|.
\end{equation}

Esta expresi\'{o}n no se toma como error experimental. M\'{a}s bien es una forma sencilla de medir qu\'{e} tan cerca pasa la trayectoria del blanco.

\section{Dise\~{n}o de la soluci\'{o}n}

El c\'{o}digo se organiz\'{o} usando programaci\'{o}n orientada a objetos. La idea fue separar responsabilidades para que cada clase tuviera un papel claro. Esto ayuda a que el programa sea m\'{a}s f\'{a}cil de leer, modificar y explicar.

\begin{table}[H]
\centering
\begin{tabular}{p{4.2cm}p{10cm}}
\toprule
\textbf{Clase} & \textbf{Funci\'{o}n dentro del proyecto} \\
\midrule
\texttt{SistemaFisico} & Clase abstracta. Define que todo sistema f\'{i}sico debe tener un m\'{e}todo para simular y otro para entregar un resumen. \\
\texttt{Proyectil} & Guarda las propiedades b\'{a}sicas del proyectil: velocidad inicial y gravedad. \\
\texttt{LanzamientoParabolico} & Calcula la trayectoria, busca el \'{a}ngulo \'{o}ptimo y genera los datos simulados. \\
\texttt{Analizador} & Muestra en consola los resultados principales del lanzamiento. \\
\texttt{Visualizador} & Crea la gr\'{a}fica y la animaci\'{o}n del proyectil. \\
\bottomrule
\end{tabular}
\caption{Clases principales del proyecto.}
\end{table}

La clase \texttt{SistemaFisico} permite cumplir la idea de abstracci\'{o}n. No representa un lanzamiento concreto, sino una estructura general. La clase \texttt{LanzamientoParabolico} hereda de ella y s\'{i} implementa el caso f\'{i}sico espec\'{i}fico.

\section{Librer\'{i}as utilizadas}

\begin{table}[H]
\centering
\begin{tabular}{p{4cm}p{10cm}}
\toprule
\textbf{Librer\'{i}a} & \textbf{Uso en el proyecto} \\
\midrule
\texttt{abc} & Permite crear la clase abstracta \texttt{SistemaFisico}. \\
\texttt{NumPy} & Se usa para senos, cosenos, conversi\'{o}n de \'{a}ngulos, arreglos de tiempo y valores m\'{a}ximos. \\
\texttt{Matplotlib} & Se usa para graficar la trayectoria, el objetivo y el punto animado. \\
\texttt{FuncAnimation} & Permite animar el movimiento del proyectil sobre la trayectoria calculada. \\
\texttt{SciPy} & Se usa para optimizar el \'{a}ngulo de lanzamiento mediante \texttt{minimize\_scalar}. \\
\bottomrule
\end{tabular}
\caption{Librer\'{i}as utilizadas en el c\'{o}digo.}
\end{table}

\section{Explicaci\'{o}n del funcionamiento del c\'{o}digo}

El c\'{o}digo empieza importando las librer\'{i}as necesarias. Despu\'{e}s se define una clase abstracta llamada \texttt{SistemaFisico}. Esta clase obliga a que las clases hijas tengan los m\'{e}todos \texttt{simular} y \texttt{resumen}. En este caso, la clase que hereda de ella es \texttt{LanzamientoParabolico}.

Luego aparece la clase \texttt{Proyectil}, que guarda la velocidad inicial y la gravedad. Esto permite separar los datos del objeto lanzado de los c\'{a}lculos de la trayectoria.

La clase \texttt{LanzamientoParabolico} es la parte central. En ella se guardan el proyectil y las coordenadas del objetivo. Tambi\'{e}n se crean variables como \texttt{angulo\_optimo}, \texttt{error\_minimo}, \texttt{t}, \texttt{x} y \texttt{y}. Al comienzo algunas de estas variables est\'{a}n en \texttt{None} porque todav\'{i}a no se ha hecho la simulaci\'{o}n.

El m\'{e}todo \texttt{posicion} implementa directamente las ecuaciones del movimiento parab\'{o}lico. Recibe un \'{a}ngulo y un tiempo, y devuelve las coordenadas \(x\) y \(y\). Esta parte es la base f\'{i}sica del programa.

El m\'{e}todo \texttt{error\_objetivo} calcula qu\'{e} tan cerca pasa la trayectoria del objetivo para un \'{a}ngulo dado. Primero calcula el tiempo en que el proyectil llega a la posici\'{o}n horizontal del objetivo. Luego eval\'{u}a la altura del proyectil en ese instante y la compara con la altura del blanco. Aunque el nombre del m\'{e}todo usa la palabra error, realmente se est\'{a} midiendo una desviaci\'{o}n vertical.

Despu\'{e}s, el m\'{e}todo \texttt{optimizar\_angulo} usa \texttt{minimize\_scalar}. Esta funci\'{o}n prueba distintos \'{a}ngulos entre 1 y 89 grados y busca cu\'{a}l minimiza la desviaci\'{o}n respecto al objetivo. Se evita usar 0 y 90 grados porque son casos extremos: en 0 grados casi no hay componente vertical, y en 90 grados no hay avance horizontal.

El m\'{e}todo \texttt{simular} se encarga de generar la trayectoria completa. Si todav\'{i}a no se ha calculado el \'{a}ngulo \'{o}ptimo, primero lo calcula. Luego crea un arreglo de tiempos con \texttt{np.linspace} y obtiene los valores de \(x\) y \(y\) usando el m\'{e}todo \texttt{posicion}.

El m\'{e}todo \texttt{resumen} devuelve los datos principales en forma de diccionario: \'{a}ngulo \'{o}ptimo, desviaci\'{o}n m\'{i}nima, altura m\'{a}xima y alcance de la trayectoria.

La clase \texttt{Analizador} toma esos resultados y los imprime en consola. Esta parte ayuda a tener una salida num\'{e}rica clara, no solamente una gr\'{a}fica. Por \'{u}ltimo, la clase \texttt{Visualizador} crea la figura, dibuja la trayectoria, marca el objetivo con una equis y anima un punto que representa el proyectil movi\'{e}ndose.

\section{C\'{o}digo utilizado}

\begin{lstlisting}[style=pythonstyle]
from abc import ABC, abstractmethod
import numpy as np
import matplotlib.pyplot as plt
from matplotlib.animation import FuncAnimation
from scipy.optimize import minimize_scalar

class SistemaFisico(ABC):
    @abstractmethod
    def simular(self):
        pass

    @abstractmethod
    def resumen(self):
        pass

class Proyectil:
    def __init__(self, velocidad_inicial, gravedad):
        self.v0 = velocidad_inicial
        self.g = gravedad

class LanzamientoParabolico(SistemaFisico):
    def __init__(self, proyectil, objetivo_x, objetivo_y):
        self.proyectil = proyectil
        self.objetivo_x = objetivo_x
        self.objetivo_y = objetivo_y
        self.angulo_optimo = None
        self.error_minimo = None
        self.t = None
        self.x = None
        self.y = None

    def posicion(self, angulo, t):
        v0 = self.proyectil.v0
        g = self.proyectil.g
        x = v0 * np.cos(angulo) * t
        y = v0 * np.sin(angulo) * t - 0.5 * g * t**2
        return x, y

    def error_objetivo(self, angulo):
        v0 = self.proyectil.v0
        t_objetivo = self.objetivo_x / (v0 * np.cos(angulo))
        _, altura = self.posicion(angulo, t_objetivo)
        return abs(altura - self.objetivo_y)

    def optimizar_angulo(self):
        resultado = minimize_scalar(
            self.error_objetivo,
            bounds=(np.radians(1), np.radians(89)),
            method="bounded"
        )
        self.angulo_optimo = resultado.x
        self.error_minimo = resultado.fun

    def simular(self):
        if self.angulo_optimo is None:
            self.optimizar_angulo()
        v0 = self.proyectil.v0
        g = self.proyectil.g
        theta = self.angulo_optimo
        t_final = 2 * v0 * np.sin(theta) / g
        self.t = np.linspace(0, t_final, 160)
        self.x, self.y = self.posicion(theta, self.t)
        return self.t, self.x, self.y

    def resumen(self):
        return {
            "angulo": np.degrees(self.angulo_optimo),
            "error": self.error_minimo,
            "altura_maxima": np.max(self.y),
            "alcance": np.max(self.x)
        }

class Analizador:
    def __init__(self, sistema):
        self.sistema = sistema

    def mostrar(self):
        datos = self.sistema.resumen()
        print("ANALISIS DEL LANZAMIENTO")
        print(f"Velocidad inicial: {self.sistema.proyectil.v0:.2f} m/s")
        print(f"Gravedad: {self.sistema.proyectil.g:.2f} m/s^2")
        print(f"Objetivo: ({self.sistema.objetivo_x:.2f}, {self.sistema.objetivo_y:.2f}) m")
        print(f"Angulo optimo: {datos['angulo']:.2f} grados")
        print(f"Error minimo: {datos['error']:.4f} m")
        print(f"Altura maxima: {datos['altura_maxima']:.2f} m")
        print(f"Alcance de la trayectoria: {datos['alcance']:.2f} m")

class Visualizador:
    def __init__(self, sistema):
        self.sistema = sistema

    def animar(self):
        t, x, y = self.sistema.simular()
        fig, ax = plt.subplots(figsize=(8, 5))

        ax.set_xlim(0, max(np.max(x), self.sistema.objetivo_x) + 10)
        ax.set_ylim(0, max(np.max(y), self.sistema.objetivo_y) + 10)
        ax.set_title("Busqueda del angulo para golpear un objetivo")
        ax.set_xlabel("x (m)")
        ax.set_ylabel("y (m)")
        ax.grid(True)

        ax.plot(x, y, "-", label="trayectoria optimizada")
        ax.scatter(self.sistema.objetivo_x, self.sistema.objetivo_y,
                   marker="x", s=110, label="objetivo")

        punto, = ax.plot([], [], "o", markersize=9)
        texto = ax.text(0.03, 0.87, "", transform=ax.transAxes)
        angulo = np.degrees(self.sistema.angulo_optimo)

        def iniciar():
            punto.set_data([], [])
            texto.set_text("")
            return punto, texto

        def actualizar(i):
            punto.set_data([x[i]], [y[i]])
            texto.set_text(
                f"t = {t[i]:.2f} s\n"
                f"angulo = {angulo:.2f} grados\n"
                f"error = {self.sistema.error_minimo:.4f} m"
            )
            return punto, texto

        animacion = FuncAnimation(
            fig, actualizar, frames=len(t),
            init_func=iniciar, interval=35, blit=True
        )

        ax.legend()
        plt.show()
        return animacion

def main():
    velocidad_inicial = 28.0
    gravedad = 9.81
    objetivo_x = 5.0
    objetivo_y = 14.0

    proyectil = Proyectil(velocidad_inicial, gravedad)
    sistema = LanzamientoParabolico(proyectil, objetivo_x, objetivo_y)

    sistema.simular()
    Analizador(sistema).mostrar()
    Visualizador(sistema).animar()

if __name__ == "__main__":
    main()
\end{lstlisting}

\section{Resultados esperados}

Al ejecutar el programa, primero se muestran resultados en consola. La salida incluye la velocidad inicial, la gravedad, la posici\'{o}n del objetivo, el \'{a}ngulo \'{o}ptimo encontrado, la desviaci\'{o}n m\'{i}nima respecto al objetivo, la altura m\'{a}xima y el alcance de la trayectoria.

Con los valores usados en el programa,

\begin{align*}
v_0 &= 28.0\ \text{m/s},\\
g &= 9.81\ \text{m/s}^2,\\
x_{\text{objetivo}} &= 5.0\ \text{m},\\
y_{\text{objetivo}} &= 14.0\ \text{m},
\end{align*}

el programa busca un \'{a}ngulo que haga que la trayectoria pase cerca de ese punto. Adem\'{a}s de la salida en consola, se abre una ventana con la gr\'{a}fica de la trayectoria. En esa gr\'{a}fica aparece la curva del lanzamiento, el objetivo marcado con una equis y un punto animado que se mueve sobre la trayectoria.

El objetivo de la animaci\'{o}n no es solo decorar el trabajo. Sirve para ver de forma directa c\'{o}mo el proyectil sube, alcanza una altura m\'{a}xima y luego baja por efecto de la gravedad. Tambi\'{e}n permite ubicar visualmente el blanco y entender si la trayectoria se aproxima a \'{e}l.

\section{An\'{a}lisis de resultados}

El resultado principal del proyecto es el \'{a}ngulo \'{o}ptimo. Este \'{a}ngulo depende de la velocidad inicial, de la gravedad y de la posici\'{o}n del objetivo. Si se cambia alguno de esos datos, el programa puede encontrar un \'{a}ngulo diferente.

La trayectoria obtenida tiene forma parab\'{o}lica, como se espera en un movimiento ideal bajo gravedad constante. En el eje horizontal el proyectil avanza de forma uniforme, mientras que en el eje vertical sube al comienzo y despu\'{e}s cae debido a la gravedad. Esto coincide con las ecuaciones usadas en el marco te\'{o}rico.

La desviaci\'{o}n m\'{i}nima respecto al objetivo permite saber qu\'{e} tan cerca pasa la trayectoria del blanco. Esta cantidad no debe confundirse con error experimental. No se est\'{a} midiendo una pr\'{a}ctica real con incertidumbre, sino usando una medida matem\'{a}tica para que el programa pueda decidir qu\'{e} \'{a}ngulo es mejor.

La altura m\'{a}xima ayuda a interpretar qu\'{e} tan alto llega el proyectil, y el alcance indica hasta qu\'{e} distancia horizontal llega antes de volver al suelo. Estos datos complementan el resultado principal, porque no solo se obtiene el \'{a}ngulo, sino tambi\'{e}n una descripci\'{o}n general del movimiento.

\section{Cumplimiento de la r\'{u}brica}

\begin{longtable}{p{5.5cm}p{9cm}}
\toprule
\textbf{Requisito} & \textbf{C\'{o}mo se cumple en el proyecto} \\
\midrule
Al menos tres clases propias & Se usan las clases \texttt{Proyectil}, \texttt{LanzamientoParabolico}, \texttt{Analizador} y \texttt{Visualizador}. \\
Clase base abstracta & Se define la clase \texttt{SistemaFisico} usando \texttt{ABC}. \\
Herencia & \texttt{LanzamientoParabolico} hereda de \texttt{SistemaFisico}. \\
Uso de NumPy & Se usa para funciones trigonom\'{e}tricas, conversi\'{o}n de \'{a}ngulos, arreglos de tiempo y valores m\'{a}ximos. \\
Uso de Matplotlib & Se usa para graficar la trayectoria, el objetivo y el punto animado. \\
Uso de SciPy & Se usa \texttt{minimize\_scalar} para buscar el \'{a}ngulo \'{o}ptimo. \\
Datos simulados & Se generan los arreglos \texttt{t}, \texttt{x} y \texttt{y}. \\
Visualizaci\'{o}n & Se muestra la trayectoria parab\'{o}lica y el objetivo. \\
Animaci\'{o}n & Se anima el proyectil usando \texttt{FuncAnimation}. \\
An\'{a}lisis de resultados & La clase \texttt{Analizador} imprime los resultados principales y estos se interpretan en el documento. \\
C\'{o}digo que compile & El programa est\'{a} organizado en un solo archivo de Python y mantiene una estructura clara. \\
Dise\~{n}o de soluci\'{o}n & El problema se dividi\'{o} en clases con responsabilidades distintas. \\
Planteamiento del problema y objetivo & Se plantea la b\'{u}squeda del \'{a}ngulo para acercarse a un objetivo fijo. \\
\bottomrule
\end{longtable}

\section{Conclusiones}

Se logr\'{o} construir una simulaci\'{o}n computacional de un lanzamiento parab\'{o}lico usando un modelo f\'{i}sico ideal. El programa calcula la trayectoria a partir de las ecuaciones de cinem\'{a}tica y permite encontrar un \'{a}ngulo de lanzamiento adecuado para aproximarse a un objetivo definido.

El uso de \texttt{SciPy} permiti\'{o} darle una parte computacional m\'{a}s fuerte al proyecto, porque el programa no se limita a graficar una curva, sino que busca num\'{e}ricamente el \'{a}ngulo que minimiza la desviaci\'{o}n respecto al blanco.

La programaci\'{o}n orientada a objetos ayud\'{o} a organizar el c\'{o}digo de una forma m\'{a}s clara. Cada clase cumple una funci\'{o}n concreta: una define la estructura general, otra guarda los datos del proyectil, otra calcula la trayectoria, otra analiza resultados y otra se encarga de la visualizaci\'{o}n.

La animaci\'{o}n tambi\'{e}n fue importante porque facilita la interpretaci\'{o}n f\'{i}sica del movimiento. Al ver el punto moverse sobre la curva se entiende mejor c\'{o}mo cambia la posici\'{o}n del proyectil con el tiempo y c\'{o}mo la gravedad afecta su trayectoria.

\end{document}
